\begin{abstract}
AI agents powered by large language models (LLMs) have shown strong capabilities in problem solving. Through combining many intelligent agents, multi-agent collaboration has emerged as a promising approach to tackle complex, multi-faceted problems that exceed the capabilities of single AI agents. However, designing the collaboration protocols and evaluating the effectiveness of these systems remains a significant challenge, especially for enterprise applications. This report addresses these challenges by presenting a comprehensive evaluation of coordination and routing capabilities in a novel multi-agent collaboration framework. We evaluate two key operational modes: (1) a coordination mode enabling complex task completion through parallel communication and payload referencing, and (2) a routing mode for efficient message forwarding between agents. We benchmark on a set of handcrafted scenarios from three enterprise domains, which are publicly released with the report. For coordination capabilities, we demonstrate the effectiveness of inter-agent communication and payload referencing mechanisms, achieving end-to-end goal success rates of 90\%.
%In routing mode, we compare LLM-based (Claude 3.0 Haiku) and Lex-based approaches. The LLM-based router achieves 90-92\% classification accuracy with minimal false switches (~3\%) and average routing overhead of 600-800ms per turn. While the Lex-based approach offers lower latency (~140ms), it demonstrates reduced accuracy (54-85\%), highlighting important trade-offs between speed and reliability. 
% Our analysis reveals several key findings: multi-agent collaboration improves goal success rates relatively by up to 70\% when utilized effectively; payload referencing enhances code-heavy tasks relatively by 23\% and routing mode significantly improves latency. These results provide insights for deploying multi-agent systems for enterprise applications and contribute to the ongoing research in scalable, efficient multi-agent collaboration frameworks.
 Our analysis yields several key findings: multi-agent collaboration enhances goal success rates by up to 70\% compared to single-agent approaches in our benchmarks; payload referencing improves performance on code-intensive tasks by 23\%; latency can be substantially reduced with a routing mechanism that selectively bypasses agent orchestration. These findings offer valuable guidance for enterprise deployments of multi-agent systems and advance the development of scalable, efficient multi-agent collaboration frameworks.

\end{abstract}